\documentclass[12pt,a4paper]{article}

\usepackage[T2A]{fontenc}
\usepackage[utf8]{inputenc}
\usepackage[russian]{babel}
\usepackage{amsmath,amssymb,amsthm,mathtools}
\usepackage{array}
\usepackage{booktabs}
\usepackage{algorithm}
\usepackage{algpseudocode}
\usepackage{geometry}
\usepackage{enumitem}
\usepackage{graphicx}
\usepackage{tikz}
\usetikzlibrary{automata,positioning,arrows.meta}

\geometry{left=25mm,right=15mm,top=20mm,bottom=20mm}

\newtheorem{definition}{Определение}
\newtheorem{theorem}{Теорема}
\newtheorem{lemma}{Лемма}
\newtheorem{proposition}{Утверждение}
\newtheorem{example}{Пример}
\newtheorem{corollary}{Заклчение}
\newtheorem{remark}{Замечание}

\DeclareMathOperator{\sgn}{sgn}
\DeclareMathOperator{\seq}{seq}
\DeclareMathOperator{\ord}{ord}

\newcommand{\eps}{\varepsilon}
\newcommand{\N}{\mathbb{N}}
\newcommand{\set}[1]{\{#1\}}
\newcommand{\powerset}[1]{2^{#1}}
\newcommand{\derives}{\Rightarrow}
\newcommand{\derivesstar}{\Rightarrow^{*}}
\newcommand{\derivesplus}{\Rightarrow^{+}}
\newcommand{\bigmid}{\mathop{\big|}}

\title{Билет №12 \\ \small Формальные языки и способы их описания. Классификация формальных грамматик. (ИТМО) \\ Формальные языки и способы их описания. Классификация формальных грамматик. Их использование в лексическом и синтаксическом анализе. (СпБГУ)}
\author{Сериков Владислав Александрович}
\date{5 мая 2026}

\begin{document}

\maketitle
\tableofcontents
\newpage

\section{Базовые понятия}

\subsection{Алфавит, цепочки, языки}

\begin{definition}
\textbf{Алфавит} --- конечное непустое множество символов, обычно обозначаемое $\Sigma$.
\end{definition}

\begin{definition}
\textbf{Цепочка} над алфавитом $\Sigma$ --- конечная последовательность символов из $\Sigma$.
Пустая цепочка обозначается $\eps$.
\end{definition}

Если $\alpha,\beta \in \Sigma^*$, то их \textbf{конкатенация} обозначается $\alpha\beta$.

\begin{definition}
\textbf{Длина цепочки} $\alpha$, обозначаемая $|\alpha|$, равна числу символов в цепочке.
Для пустой цепочки $|\eps|=0$.
\end{definition}

Степень цепочки определяется рекурсивно:
\[
\alpha^0=\eps,\qquad \alpha^{n+1}=\alpha^n\alpha.
\]

\begin{definition}
Для алфавита $\Sigma$:
\[
\Sigma^0=\{\eps\}, \qquad
\Sigma^n=\{\text{все цепочки длины } n \text{ над } \Sigma\},
\]
\[
\Sigma^+=\bigcup_{n\ge 1}\Sigma^n,\qquad
\Sigma^*=\bigcup_{n\ge 0}\Sigma^n.
\]
Множество $\Sigma^*$ называется \textbf{итерацией алфавита}, а $\Sigma^+$ --- \textbf{усечённой итерацией}.
\end{definition}

\begin{definition}
\textbf{Формальный язык} над алфавитом $\Sigma$ --- любое множество цепочек над $\Sigma$, то есть любое подмножество:
\[
L\subseteq \Sigma^*.
\]
\end{definition}

Таким образом, формальный язык может быть конечным, бесконечным или пустым.

\begin{example}
Пусть $\Sigma=\{a,b\}$. Тогда:
\[
L_1=\{a^n b^n\mid n\ge 1\}
\]
--- язык цепочек, состоящих из некоторого положительного числа букв $a$, за которыми следует такое же число букв $b$:
\[
L_1=\{ab,aabb,aaabbb,\ldots\}.
\]
\end{example}

\subsection{Основные способы задания формальных языков}

Формальные языки задаются несколькими основными способами.

\begin{enumerate}
    \item \textbf{Перечислением элементов}, если язык конечен:
    \[
    L=\{ab,ba,aa\}.
    \]

    \item \textbf{С помощью характеристического условия}:
    \[
    L=\{\omega\in\{0,1\}^* \mid \omega \text{ содержит чётное число единиц}\}.
    \]

    \item \textbf{С помощью порождающей грамматики}.
    Грамматика задаёт правила построения цепочек языка.

    \item \textbf{С помощью распознавателя}.
    Распознаватель получает цепочку на вход и отвечает, принадлежит ли она языку.

    \item \textbf{С помощью регулярного выражения}, если язык регулярный:
    \[
    (a|b)^*abb.
    \]

    \item \textbf{С помощью автоматов}.
    Регулярные языки задаются конечными автоматами, контекстно-свободные --- автоматами с магазинной памятью.
\end{enumerate}

\section{Порождающие грамматики}

\subsection{Определение грамматики}

\begin{definition}
\textbf{Порождающая грамматика} --- четвёрка
\[
G=(V_T,V_N,P,S),
\]
где:
\begin{itemize}
    \item $V_T$ --- конечное множество терминальных символов;
    \item $V_N$ --- конечное множество нетерминальных символов;
    \item $V_T\cap V_N=\varnothing$;
    \item $P$ --- конечное множество правил вывода;
    \item $S\in V_N$ --- начальный символ грамматики.
\end{itemize}
\end{definition}

Полный словарь грамматики:
\[
V=V_T\cup V_N.
\]

Правило вывода имеет вид
\[
\alpha\to\beta,
\]
где $\alpha,\beta$ --- цепочки над $V$, причём левая часть $\alpha$ не должна быть пустой.

\begin{definition}
Если цепочка $\gamma$ получается из цепочки $\delta$ заменой в $\delta$ некоторой подцепочки, совпадающей с левой частью правила, на правую часть этого правила, то говорят, что $\delta$ \textbf{непосредственно порождает} $\gamma$:
\[
\delta\derives \gamma.
\]
\end{definition}

Если существует конечная последовательность непосредственных выводов:
\[
\alpha=\alpha_0\derives \alpha_1\derives \ldots \derives \alpha_n=\beta,
\]
то пишут:
\[
\alpha\derivesstar \beta.
\]

\begin{definition}
\textbf{Язык, порождаемый грамматикой} $G$, определяется так:
\[
L(G)=\{\omega\in V_T^*\mid S\derivesstar \omega\}.
\]
\end{definition}

\begin{definition}
Цепочка $\alpha\in V^*$ называется \textbf{сентенциальной формой}, если
\[
S\derivesstar \alpha.
\]
Если при этом $\alpha\in V_T^*$, то $\alpha$ называется \textbf{предложением языка}.
\end{definition}

\subsection{Левосторонний и правосторонний вывод}

\begin{definition}
Вывод называется \textbf{левосторонним}, если на каждом шаге раскрывается самый левый нетерминал сентенциальной формы.
\end{definition}

\begin{definition}
Вывод называется \textbf{правосторонним}, если на каждом шаге раскрывается самый правый нетерминал сентенциальной формы.
\end{definition}

\begin{example}
Рассмотрим грамматику:
\[
S\to aSb\mid ab.
\]
Она порождает язык
\[
L=\{a^n b^n\mid n\ge 1\}.
\]

Вывод цепочки $aaabbb$:
\[
S\derives aSb\derives aaSbb\derives aaabbb.
\]
\end{example}

\section{Классификация формальных грамматик по Хомскому}

Классификация Хомского делит грамматики на четыре типа в зависимости от формы правил.

\subsection{Грамматики типа 0: неограниченные грамматики}

\begin{definition}
Грамматика называется \textbf{грамматикой типа 0}, или \textbf{неограниченной грамматикой}, если её правила имеют вид
\[
\alpha\to\beta,
\]
где
\[
\alpha\in V^+,\qquad \beta\in V^*.
\]
\end{definition}

Ограничение только одно: левая часть правила непуста.

Языки типа 0 называются \textbf{рекурсивно перечислимыми}. Их распознают машины Тьюринга.

\subsection{Грамматики типа 1: контекстно-зависимые грамматики}

\begin{definition}
Грамматика называется \textbf{контекстно-зависимой}, если её правила имеют вид
\[
\alpha A\beta \to \alpha \gamma \beta,
\]
где
\[
A\in V_N,\qquad \alpha,\beta\in V^*,\qquad \gamma\in V^+.
\]
\end{definition}

Иначе говоря, нетерминал $A$ можно заменить на $\gamma$ только в контексте $\alpha$ и $\beta$.

Эквивалентная форма определения: грамматика типа 1 является \textbf{неукорачивающей}, то есть для каждого правила
\[
\alpha\to\beta
\]
выполнено
\[
|\alpha|\le |\beta|.
\]

Исключение иногда допускается для правила
\[
S\to \eps,
\]
если начальный символ $S$ не встречается в правых частях правил.

\begin{example}
Язык
\[
L=\{a^n b^n c^n\mid n\ge 1\}
\]
является контекстно-зависимым, но не является контекстно-свободным.
\end{example}

\subsection{Грамматики типа 2: контекстно-свободные грамматики}

\begin{definition}
Грамматика называется \textbf{контекстно-свободной}, если каждое её правило имеет вид
\[
A\to \alpha,
\]
где
\[
A\in V_N,\qquad \alpha\in V^*.
\]
\end{definition}

Левая часть правила состоит ровно из одного нетерминала.

Контекстно-свободные языки распознаются автоматами с магазинной памятью.

\begin{example}
Грамматика
\[
S\to aSb\mid \eps
\]
порождает язык
\[
L=\{a^n b^n\mid n\ge 0\}.
\]
Это контекстно-свободный язык.
\end{example}

\subsection{Грамматики типа 3: регулярные грамматики}

\begin{definition}
Грамматика называется \textbf{праволинейной регулярной}, если каждое правило имеет вид
\[
A\to aB
\]
или
\[
A\to a,
\]
где
\[
A,B\in V_N,\qquad a\in V_T.
\]
\end{definition}

\begin{definition}
Грамматика называется \textbf{леволинейной регулярной}, если каждое правило имеет вид
\[
A\to Ba
\]
или
\[
A\to a.
\]
\end{definition}

Регулярные языки распознаются конечными автоматами и задаются регулярными выражениями.

\begin{example}
Грамматика
\[
S\to aS\mid bA,\qquad A\to bA\mid b
\]
регулярна. Она порождает язык:
\[
L=\{a^n b^m\mid n\ge 0,\ m\ge 2\}.
\]
\end{example}

\subsection{Иерархия Хомского}

Классы языков образуют строгую иерархию:
\[
\text{Регулярные}
\subset
\text{Контекстно-свободные}
\subset
\text{Контекстно-зависимые}
\subset
\text{Рекурсивно перечислимые}.
\]

\[
\mathcal{L}_3\subset \mathcal{L}_2\subset \mathcal{L}_1\subset \mathcal{L}_0.
\]

\begin{center}
\begin{tabular}{>{\centering\arraybackslash}m{2cm}
                >{\centering\arraybackslash}m{4cm}
                >{\centering\arraybackslash}m{5cm}}
\toprule
Тип & Класс грамматик & Распознаватель \\
\midrule
0 & Неограниченные & Машина Тьюринга \\
1 & Контекстно-зависимые & Линейно-ограниченный автомат \\
2 & Контекстно-свободные & Автомат с магазинной памятью \\
3 & Регулярные & Конечный автомат \\
\bottomrule
\end{tabular}
\end{center}

\section{Регулярные языки}

\subsection{Регулярные выражения}

\begin{definition}
Регулярные выражения над алфавитом $\Sigma$ определяются рекурсивно:
\begin{enumerate}
    \item $\varnothing$ --- регулярное выражение;
    \item $\eps$ --- регулярное выражение;
    \item если $a\in\Sigma$, то $a$ --- регулярное выражение;
    \item если $r$ и $s$ --- регулярные выражения, то
    \[
    (r|s),\qquad (rs),\qquad (r^*)
    \]
    также регулярные выражения.
\end{enumerate}
\end{definition}

Операции:
\begin{itemize}
    \item $r|s$ --- объединение;
    \item $rs$ --- конкатенация;
    \item $r^*$ --- итерация Клини.
\end{itemize}

\begin{example}
Регулярное выражение
\[
(a|b)^*abb
\]
задаёт язык всех цепочек над $\{a,b\}$, заканчивающихся на $abb$.
\end{example}

\subsection{Конечные автоматы}

\begin{definition}
\textbf{Детерминированный конечный автомат} --- пятёрка
\[
A=(\Sigma,Q,q_0,\delta,F),
\]
где:
\begin{itemize}
    \item $\Sigma$ --- входной алфавит;
    \item $Q$ --- конечное множество состояний;
    \item $q_0\in Q$ --- начальное состояние;
    \item $\delta:Q\times \Sigma\to Q$ --- функция переходов;
    \item $F\subseteq Q$ --- множество допускающих состояний.
\end{itemize}
\end{definition}

Автомат читает входную цепочку слева направо. Если после чтения всей цепочки он оказывается в состоянии из $F$, цепочка допускается.

\begin{definition}
Язык, допускаемый автоматом $A$, равен
\[
L(A)=\{\omega\in\Sigma^*\mid \delta^*(q_0,\omega)\in F\},
\]
где $\delta^*$ --- расширенная функция переходов.
\end{definition}

Расширенная функция переходов:
\[
\delta^*(q,\eps)=q,
\]
\[
\delta^*(q,a\omega)=\delta^*(\delta(q,a),\omega).
\]

\subsection{Недетерминированные конечные автоматы}

\begin{definition}
\textbf{Недетерминированный конечный автомат} --- пятёрка
\[
A=(\Sigma,Q,q_0,\delta,F),
\]
где
\[
\delta:Q\times\Sigma\to 2^Q.
\]
\end{definition}

В отличие от ДКА, из одного состояния по одному символу может быть несколько переходов.

Также допускаются $\eps$-переходы:
\[
\delta:Q\times(\Sigma\cup\{\eps\})\to 2^Q.
\]

\begin{theorem}[эквивалентность ДКА и НКА]
Для каждого недетерминированного конечного автомата существует эквивалентный детерминированный конечный автомат, допускающий тот же язык.
\end{theorem}

\begin{proof}
Пусть дан НКА
\[
N=(\Sigma,Q,q_0,\delta,F).
\]
Построим ДКА
\[
D=(\Sigma,Q',q_0',\delta',F').
\]

Идея построения: одно состояние ДКА соответствует множеству возможных состояний НКА.

Положим:
\[
Q'=2^Q.
\]
Начальное состояние:
\[
q_0'=\{q_0\}.
\]
Если в НКА есть $\eps$-переходы, то вместо $\{q_0\}$ берётся $\eps$-замыкание начального состояния:
\[
q_0'=\eps\text{-closure}(q_0).
\]

Функция переходов задаётся так:
\[
\delta'(R,a)=\bigcup_{q\in R}\delta(q,a),
\]
где $R\subseteq Q$.

Если есть $\eps$-переходы, то:
\[
\delta'(R,a)=\eps\text{-closure}\left(\bigcup_{q\in R}\delta(q,a)\right).
\]

Допускающие состояния ДКА:
\[
F'=\{R\subseteq Q\mid R\cap F\ne\varnothing\}.
\]

Докажем корректность построения.

После чтения любой цепочки $\omega$ состояние ДКА равно множеству всех состояний, в которых мог бы находиться НКА после чтения $\omega$.

Доказывается индукцией по длине $\omega$.

\textbf{База.} Для $\omega=\eps$ ДКА находится в $\{q_0\}$ или в $\eps$-замыкании $q_0$, что совпадает с множеством состояний НКА до чтения входа.

\textbf{Переход.} Пусть утверждение верно для $\omega$. Рассмотрим $\omega a$.
После чтения $\omega$ ДКА находится во множестве $R$, равном множеству всех возможных состояний НКА. При чтении символа $a$ НКА может перейти в любое состояние из
\[
\bigcup_{q\in R}\delta(q,a).
\]
Именно это множество и задаёт переход ДКА. При наличии $\eps$-переходов дополнительно берётся $\eps$-замыкание.

Следовательно, после чтения любой цепочки ДКА точно моделирует все возможные вычисления НКА.

Цепочка допускается НКА тогда и только тогда, когда среди возможных состояний после её чтения есть хотя бы одно допускающее состояние. По определению $F'$ это эквивалентно тому, что соответствующее состояние ДКА принадлежит $F'$.

Значит,
\[
L(D)=L(N).
\]
\end{proof}

\subsection{Теорема Клини}

\begin{theorem}[Клини]
Язык является регулярным тогда и только тогда, когда он допускается некоторым конечным автоматом.
\end{theorem}

\begin{proof}
Докажем обе стороны.

\textbf{1. Если язык задан регулярным выражением, то он допускается конечным автоматом.}

Построение ведётся индукцией по структуре регулярного выражения.

\begin{enumerate}
    \item Для $\varnothing$ строится автомат без допускающих путей.
    \item Для $\eps$ строится автомат, в котором начальное состояние является допускающим.
    \item Для символа $a$ строится автомат:
    \[
    q_0 \xrightarrow{a} q_1,
    \]
    где $q_1$ допускающее.
\end{enumerate}

Для составных выражений:
\begin{itemize}
    \item для $r|s$ строится новый начальный узел с $\eps$-переходами к автоматам для $r$ и $s$;
    \item для $rs$ допускающие состояния автомата для $r$ соединяются $\eps$-переходами с начальным состоянием автомата для $s$;
    \item для $r^*$ добавляется возможность пройти автомат для $r$ любое число раз, включая ноль.
\end{itemize}

Так получается НКА с $\eps$-переходами. По теореме об эквивалентности НКА и ДКА он может быть преобразован в ДКА.

Следовательно, всякий язык, заданный регулярным выражением, допускается конечным автоматом.

\textbf{2. Если язык допускается конечным автоматом, то он задаётся регулярным выражением.}

Пусть дан конечный автомат
\[
A=(\Sigma,Q,q_0,\delta,F).
\]
Состояния можно перенумеровать:
\[
Q=\{1,2,\ldots,n\}.
\]

Введём регулярное выражение $R_{ij}^{(k)}$, задающее множество всех цепочек, переводящих автомат из состояния $i$ в состояние $j$ и проходящих в качестве промежуточных только через состояния из множества $\{1,\ldots,k\}$.

Для $k=0$ выражение $R_{ij}^{(0)}$ задаётся непосредственно переходами:
\[
R_{ij}^{(0)}=
\begin{cases}
\text{сумма всех символов } a,\text{ таких что } \delta(i,a)=j, & i\ne j,\\
\eps \mid \text{сумма всех символов } a,\text{ таких что } \delta(i,a)=i, & i=j.
\end{cases}
\]

Рекуррентно:
\[
R_{ij}^{(k)}
=
R_{ij}^{(k-1)}
\mid
R_{ik}^{(k-1)}
\left(R_{kk}^{(k-1)}\right)^*
R_{kj}^{(k-1)}.
\]

Это выражение учитывает два случая:
\begin{enumerate}
    \item путь из $i$ в $j$ вообще не проходит через состояние $k$;
    \item путь проходит через $k$: сначала из $i$ в $k$, затем некоторое число циклов в $k$, затем из $k$ в $j$.
\end{enumerate}

После построения $R_{ij}^{(n)}$ получаем регулярное выражение для языка автомата:
\[
R=\bigmid_{f\in F} R_{q_0 f}^{(n)}.
\]

Значит, язык любого конечного автомата задаётся регулярным выражением.

Обе стороны доказаны, поэтому классы регулярных выражений и конечных автоматов совпадают.
\end{proof}

\section{Алгоритмы для регулярных языков}

\subsection{Построение НКА по регулярному выражению}

Используется конструкция Томпсона.

\subsubsection*{Базовые случаи}

\[
\eps:
\quad
q_0\xrightarrow{\eps}q_1.
\]

\[
a:
\quad
q_0\xrightarrow{a}q_1.
\]

\subsubsection*{Объединение}

Для выражения $r|s$:

\begin{center}
\begin{tikzpicture}[shorten >=1pt,node distance=2cm,on grid,auto]
\node[state] (q0) {$q_0$};
\node[state] (r) [above right=of q0] {$N_r$};
\node[state] (s) [below right=of q0] {$N_s$};
\node[state] (qf) [right=4cm of q0] {$q_f$};
\path[->]
(q0) edge node {$\eps$} (r)
(q0) edge node {$\eps$} (s)
(r) edge node {$\eps$} (qf)
(s) edge node {$\eps$} (qf);
\end{tikzpicture}
\end{center}

\subsubsection*{Конкатенация}

Для выражения $rs$ допускающие состояния автомата $N_r$ соединяются $\eps$-переходами с начальным состоянием $N_s$.

\subsubsection*{Итерация}

Для выражения $r^*$ добавляется новый начальный и новый допускающий узел:
\[
q_0 \xrightarrow{\eps} N_r,\qquad
q_0\xrightarrow{\eps}q_f,
\]
а из конца $N_r$ добавляются переходы:
\[
N_r \xrightarrow{\eps} q_f,\qquad
N_r \xrightarrow{\eps} \text{начало } N_r.
\]

\subsection{Построение ДКА по НКА}

\textbf{Вход:} НКА $N=(\Sigma,Q,q_0,\delta,F)$.

\textbf{Выход:} ДКА $D$, такой что $L(D)=L(N)$.

\subsubsection*{Шаги алгоритма}

\begin{enumerate}
    \item Начальным состоянием ДКА сделать множество:
    \[
    \eps\text{-closure}(q_0).
    \]

    \item Для каждого ещё не обработанного состояния ДКА $T\subseteq Q$ и каждого символа $a\in\Sigma$ вычислить:
    \[
    U=\eps\text{-closure}(\operatorname{move}(T,a)).
    \]

    \item Если множество $U$ ещё не было добавлено как состояние ДКА, добавить его.

    \item Добавить переход:
    \[
    T\xrightarrow{a}U.
    \]

    \item Повторять, пока не останется необработанных состояний.

    \item Допускающими состояниями ДКА объявить все множества $T$, такие что:
    \[
    T\cap F\ne\varnothing.
    \]
\end{enumerate}

\subsection{Минимизация ДКА}

\begin{definition}
Два состояния $p$ и $q$ ДКА называются \textbf{эквивалентными}, если для любой цепочки $\omega\in\Sigma^*$ выполнено:
\[
\delta^*(p,\omega)\in F
\Longleftrightarrow
\delta^*(q,\omega)\in F.
\]
\end{definition}

\subsubsection*{Алгоритм минимизации}

\begin{enumerate}
    \item Удалить недостижимые состояния.
    \item Разбить множество состояний на два класса:
    \[
    F,\qquad Q\setminus F.
    \]
    \item Последовательно уточнять разбиение: два состояния остаются в одном классе, если по каждому символу алфавита они переходят в состояния одного и того же класса.
    \item Повторять уточнение, пока разбиение не перестанет изменяться.
    \item Каждый класс эквивалентности заменить одним состоянием.
\end{enumerate}

\begin{example}
Пусть ДКА распознаёт язык цепочек над $\{a,b\}$, содержащих подцепочку $bb$.

Минимальный автомат имеет три состояния:
\begin{itemize}
    \item $q_0$ --- пока не видели $b$ как возможное начало $bb$;
    \item $q_1$ --- последний символ был $b$;
    \item $q_2$ --- уже встретили $bb$.
\end{itemize}

Переходы:
\[
\delta(q_0,a)=q_0,\quad \delta(q_0,b)=q_1,
\]
\[
\delta(q_1,a)=q_0,\quad \delta(q_1,b)=q_2,
\]
\[
\delta(q_2,a)=q_2,\quad \delta(q_2,b)=q_2.
\]
Допускающее состояние:
\[
F=\{q_2\}.
\]
\end{example}

\section{Лемма о накачке для регулярных языков}

\begin{lemma}[о накачке для регулярных языков]
Если $L$ --- регулярный язык, то существует число $p\ge 1$, такое что любая цепочка $\omega\in L$ длины не меньше $p$ может быть представлена в виде
\[
\omega=xyz,
\]
где:
\[
|xy|\le p,\qquad |y|\ge 1,
\]
и для любого $i\ge 0$:
\[
xy^i z\in L.
\]
\end{lemma}

\begin{proof}
Так как $L$ регулярный, существует ДКА
\[
A=(\Sigma,Q,q_0,\delta,F),
\]
такой что $L=L(A)$.

Положим:
\[
p=|Q|.
\]

Возьмём любую цепочку
\[
\omega=a_1a_2\ldots a_n\in L,
\]
где $n\ge p$.

При чтении первых $p$ символов автомат проходит последовательность из $p+1$ состояний:
\[
q_0,q_1,\ldots,q_p.
\]
Так как состояний всего $p$, по принципу Дирихле среди этих $p+1$ состояний найдутся два одинаковых:
\[
q_i=q_j,\qquad 0\le i<j\le p.
\]

Разложим цепочку:
\[
x=a_1\ldots a_i,
\]
\[
y=a_{i+1}\ldots a_j,
\]
\[
z=a_{j+1}\ldots a_n.
\]

Тогда:
\[
|xy|=j\le p,
\]
\[
|y|=j-i\ge 1.
\]

Так как $q_i=q_j$, чтение подцепочки $y$ переводит автомат из состояния $q_i$ снова в $q_i$. Значит, этот цикл можно пройти любое число раз.

Следовательно, для любого $i\ge 0$ цепочка
\[
xy^i z
\]
переведёт автомат из начального состояния в то же допускающее состояние, что и исходная цепочка $\omega$.

Значит:
\[
xy^i z\in L.
\]
\end{proof}

\begin{example}
Докажем, что
\[
L=\{a^n b^n\mid n\ge 1\}
\]
не является регулярным.

Предположим противное: $L$ регулярный. Тогда для него существует число накачки $p$.

Возьмём цепочку:
\[
\omega=a^p b^p.
\]
По лемме:
\[
\omega=xyz,\qquad |xy|\le p,\qquad |y|\ge 1.
\]

Так как первые $p$ символов цепочки --- это только $a$, подцепочка $y$ состоит только из букв $a$:
\[
y=a^k,\qquad k\ge 1.
\]

Накачаем вниз, взяв $i=0$:
\[
xz=a^{p-k}b^p.
\]

В этой цепочке число букв $a$ меньше числа букв $b$, поэтому:
\[
xz\notin L.
\]

Противоречие с леммой о накачке. Следовательно, язык не регулярный.
\end{example}

\section{Свойства регулярных языков}

Класс регулярных языков замкнут относительно операций:
\begin{itemize}
    \item объединения;
    \item пересечения;
    \item дополнения;
    \item разности;
    \item конкатенации;
    \item итерации Клини;
    \item обращения;
    \item гомоморфизма.
\end{itemize}

Для регулярных языков разрешимы следующие проблемы:
\begin{itemize}
    \item принадлежность цепочки языку;
    \item пустота языка;
    \item конечность языка;
    \item эквивалентность двух регулярных языков.
\end{itemize}

\section{Контекстно-свободные языки}

\subsection{Синтаксические деревья}

\begin{definition}
\textbf{Синтаксическое дерево} для КС-грамматики --- дерево, удовлетворяющее условиям:
\begin{itemize}
    \item корень помечен начальным символом $S$;
    \item внутренние вершины помечены нетерминалами;
    \item если вершина помечена $A$, а её дети слева направо помечены $X_1,\ldots,X_k$, то в грамматике есть правило
    \[
    A\to X_1X_2\ldots X_k;
    \]
    \item листья дерева, прочитанные слева направо, образуют выводимую цепочку.
\end{itemize}
\end{definition}

\begin{example}
Для грамматики
\[
S\to aSb\mid ab
\]
цепочка $aaabbb$ имеет дерево, соответствующее выводу:
\[
S\derives aSb\derives aaSbb\derives aaabbb.
\]
\end{example}

\subsection{Неоднозначность}

\begin{definition}
КС-грамматика называется \textbf{неоднозначной}, если существует цепочка языка, имеющая два различных синтаксических дерева.
\end{definition}

Эквивалентно: существует цепочка, имеющая два различных левосторонних или два различных правосторонних вывода.

\begin{example}
Грамматика
\[
E\to E+E\mid E*E\mid a
\]
неоднозначна.

Цепочка
\[
a+a*a
\]
может быть разобрана как
\[
(a+a)*a
\]
или как
\[
a+(a*a).
\]
\end{example}

\section{Автоматы с магазинной памятью}

\begin{definition}
\textbf{Автомат с магазинной памятью} --- семёрка
\[
M=(\Sigma,Q,q_0,F,\delta,\Gamma,Z_0),
\]
где:
\begin{itemize}
    \item $\Sigma$ --- входной алфавит;
    \item $Q$ --- конечное множество состояний;
    \item $q_0\in Q$ --- начальное состояние;
    \item $F\subseteq Q$ --- множество допускающих состояний;
    \item $\Gamma$ --- алфавит стека;
    \item $Z_0\in\Gamma$ --- начальный символ стека;
    \item
    \[
    \delta:Q\times(\Sigma\cup\{\eps\})\times\Gamma
    \to 2^{Q\times\Gamma^*}
    \]
    --- функция переходов.
\end{itemize}
\end{definition}

Конфигурация МП-автомата имеет вид:
\[
(q,\omega,\gamma),
\]
где:
\begin{itemize}
    \item $q$ --- текущее состояние;
    \item $\omega$ --- непрочитанная часть входа;
    \item $\gamma$ --- содержимое стека.
\end{itemize}

\subsection{Допуск по конечному состоянию и по пустому стеку}

Цепочка допускается \textbf{по конечному состоянию}, если:
\[
(q_0,\omega,Z_0)\vdash^*(q,\eps,\gamma),
\]
где
\[
q\in F.
\]

Цепочка допускается \textbf{по пустому стеку}, если:
\[
(q_0,\omega,Z_0)\vdash^*(q,\eps,\eps).
\]

Для недетерминированных МП-автоматов оба способа допуска задают один и тот же класс языков --- контекстно-свободные языки.

\begin{theorem}
Язык является контекстно-свободным тогда и только тогда, когда он допускается некоторым недетерминированным автоматом с магазинной памятью.
\end{theorem}

\begin{proof}
Докажем обе стороны.

\textbf{1. По КС-грамматике построим МП-автомат.}

Пусть дана КС-грамматика
\[
G=(V_T,V_N,P,S).
\]
Построим МП-автомат, который моделирует левосторонний вывод.

Идея:
\begin{itemize}
    \item в стеке хранится текущая сентенциальная форма, которую ещё нужно сопоставить со входом;
    \item если на вершине стека терминал, он должен совпасть с текущим входным символом;
    \item если на вершине стека нетерминал, автомат недетерминированно заменяет его правой частью одного из правил.
\end{itemize}

Автомат имеет одно состояние $q$.
Начальный символ стека --- $S$.

Для каждого правила
\[
A\to \alpha
\]
добавим переход:
\[
\delta(q,\eps,A)\ni(q,\alpha).
\]

Для каждого терминала $a\in V_T$ добавим переход:
\[
\delta(q,a,a)\ni(q,\eps).
\]

Автомат допускает по пустому стеку.

Если грамматика выводит цепочку $\omega$, то автомат может выбирать ровно те правила, которые используются в левостороннем выводе, и затем сопоставить получившиеся терминалы со входом.

Обратно, если автомат допустил цепочку, то последовательность замен нетерминалов в стеке соответствует корректному выводу в грамматике.

Значит, язык грамматики допускается построенным МП-автоматом.

\textbf{2. По МП-автомату построим КС-грамматику.}

Пусть дан МП-автомат
\[
M=(\Sigma,Q,q_0,F,\delta,\Gamma,Z_0).
\]
Без ограничения общности можно считать, что автомат допускает по пустому стеку и каждый переход либо кладёт в стек не более двух символов, либо снимает один символ.

Построим грамматику, нетерминалы которой имеют вид:
\[
[pAq],
\]
где $p,q\in Q$, $A\in\Gamma$.

Смысл нетерминала $[pAq]$:
он порождает все цепочки, которые переводят автомат из состояния $p$ со стековым символом $A$ на вершине в состояние $q$, удалив этот символ $A$.

Если автомат имеет переход, снимающий $A$:
\[
(p,a,A)\to(q,\eps),
\]
то добавляем правило:
\[
[pAq]\to a,
\]
где $a$ может быть и $\eps$.

Если автомат имеет переход:
\[
(p,a,A)\to(r,BC),
\]
то для любых состояний $s,q\in Q$ добавляем правило:
\[
[pAq]\to a[rBs][sCq].
\]

Начальный символ грамматики:
\[
S_G\to [q_0 Z_0 q]
\]
для всех $q\in Q$.

Такая грамматика порождает ровно те цепочки, которые переводят МП-автомат из начальной конфигурации в конфигурацию с пустым стеком.

Следовательно, язык любого МП-автомата является контекстно-свободным.

Итак, классы КС-языков и языков, допускаемых недетерминированными МП-автоматами, совпадают.
\end{proof}

\section{Упрощение контекстно-свободных грамматик}

\subsection{Бесполезные символы}

\begin{definition}
Нетерминал $A$ называется \textbf{производящим}, если существует цепочка терминалов $\omega\in V_T^*$, такая что:
\[
A\derivesstar \omega.
\]
\end{definition}

\begin{definition}
Символ $X\in V_T\cup V_N$ называется \textbf{достижимым}, если существуют $\alpha,\beta\in V^*$, такие что:
\[
S\derivesstar \alpha X\beta.
\]
\end{definition}

Бесполезные символы:
\begin{itemize}
    \item непроизводящие;
    \item недостижимые.
\end{itemize}

\subsubsection*{Алгоритм удаления непроизводящих символов}

\begin{enumerate}
    \item Изначально множество производящих нетерминалов пусто.
    \item Если есть правило
    \[
    A\to \alpha,
    \]
    где $\alpha$ состоит только из терминалов и уже найденных производящих нетерминалов, добавить $A$.
    \item Повторять, пока множество не перестанет изменяться.
    \item Удалить все правила, содержащие непроизводящие нетерминалы.
\end{enumerate}

\subsubsection*{Алгоритм удаления недостижимых символов}

\begin{enumerate}
    \item Изначально достижим только $S$.
    \item Если $A$ достижим и есть правило
    \[
    A\to X_1\ldots X_k,
    \]
    то все $X_i$ достижимы.
    \item Повторять до стабилизации.
    \item Удалить все правила, содержащие недостижимые символы.
\end{enumerate}

\subsection{Удаление $\eps$-правил}

\begin{definition}
Нетерминал $A$ называется \textbf{nullable}, если:
\[
A\derivesstar \eps.
\]
\end{definition}

\subsubsection*{Алгоритм}

\begin{enumerate}
    \item Найти множество всех nullable-нетерминалов.
    \item Для каждого правила
    \[
    A\to X_1X_2\ldots X_k
    \]
    добавить все варианты, полученные удалением некоторых nullable-символов.
    \item Удалить исходные $\eps$-правила.
    \item Если $\eps\in L(G)$, сохранить правило
    \[
    S\to\eps
    \]
    для нового начального символа.
\end{enumerate}

\subsection{Удаление цепных правил}

\begin{definition}
Правило вида
\[
A\to B,
\]
где $A,B\in V_N$, называется \textbf{цепным}.
\end{definition}

\subsubsection*{Алгоритм}

\begin{enumerate}
    \item Для каждого нетерминала $A$ найти множество:
    \[
    Chain(A)=\{B\mid A\derivesstar B \text{ только цепными правилами}\}.
    \]
    \item Если $B\in Chain(A)$ и есть нецепное правило
    \[
    B\to \alpha,
    \]
    добавить правило
    \[
    A\to \alpha.
    \]
    \item Удалить все цепные правила.
\end{enumerate}

\section{Нормальные формы КС-грамматик}

\subsection{Нормальная форма Хомского}

\begin{definition}
КС-грамматика находится в \textbf{нормальной форме Хомского}, если все её правила имеют один из видов:
\[
A\to BC,
\]
\[
A\to a,
\]
где
\[
A,B,C\in V_N,\qquad a\in V_T.
\]
Дополнительно может быть разрешено правило
\[
S\to\eps,
\]
если язык содержит пустую цепочку.
\end{definition}

\begin{theorem}
Любая КС-грамматика, не порождающая только пустую цепочку, может быть преобразована в эквивалентную грамматику в нормальной форме Хомского.
\end{theorem}

\begin{proof}
Пусть дана КС-грамматика $G$.

Преобразование выполняется по шагам.

\textbf{Шаг 1.} Удалим бесполезные символы.

Это не меняет язык, так как бесполезные символы не участвуют ни в одном выводе терминальной цепочки.

\textbf{Шаг 2.} Удалим $\eps$-правила.

Для каждого nullable-нетерминала добавляются все варианты правил, соответствующие его возможному исчезновению. Поэтому все выводы, в которых использовались $\eps$-правила, сохраняются без явного использования этих правил.

Если $\eps$ принадлежит языку, добавляется специальное правило $S\to\eps$.

\textbf{Шаг 3.} Удалим цепные правила.

Если
\[
A\derivesstar B
\]
цепными правилами и
\[
B\to \alpha
\]
--- нецепное правило, то добавляется правило:
\[
A\to\alpha.
\]
Тем самым любой вывод через цепочку нетерминалов заменяется одним шагом.

\textbf{Шаг 4.} Приведём длинные правые части к бинарному виду.

Если есть правило:
\[
A\to X_1X_2\ldots X_k,\qquad k\ge 3,
\]
заменим его правилами:
\[
A\to X_1Y_1,
\]
\[
Y_1\to X_2Y_2,
\]
\[
\ldots
\]
\[
Y_{k-2}\to X_{k-1}X_k.
\]

\textbf{Шаг 5.} Уберём терминалы из длинных правых частей.

Если в правиле длины 2 встречается терминал $a$, вводим новый нетерминал $A_a$ и правило:
\[
A_a\to a.
\]
Затем заменяем $a$ на $A_a$ в правых частях длины 2.

После этих шагов все правила имеют вид:
\[
A\to BC
\]
или
\[
A\to a.
\]

Каждый шаг сохраняет порождаемый язык, значит, итоговая грамматика эквивалентна исходной.
\end{proof}

\subsection{Нормальная форма Грейбах}

\begin{definition}
КС-грамматика находится в \textbf{нормальной форме Грейбах}, если каждое правило имеет вид:
\[
A\to a\alpha,
\]
где
\[
A\in V_N,\qquad a\in V_T,\qquad \alpha\in V_N^*.
\]
\end{definition}

Главное свойство: при каждом шаге вывода появляется ровно один терминал слева.

\section{Лемма о накачке для КС-языков}

\begin{lemma}[о накачке для контекстно-свободных языков]
Если $L$ --- контекстно-свободный язык, то существует число $p\ge 1$, такое что любая цепочка $\omega\in L$ длины не меньше $p$ может быть представлена в виде
\[
\omega=uvxyz,
\]
где:
\[
|vxy|\le p,
\]
\[
|vy|\ge 1,
\]
и для любого $i\ge 0$:
\[
uv^i x y^i z\in L.
\]
\end{lemma}

\begin{proof}
Так как $L$ контекстно-свободный, существует КС-грамматика $G$ в нормальной форме Хомского, порождающая $L$.

Пусть в грамматике $m$ нетерминалов.

Выберем
\[
p=2^m.
\]

Рассмотрим цепочку $\omega\in L$, где $|\omega|\ge p$.

Так как грамматика находится в нормальной форме Хомского, каждое внутреннее ветвление в синтаксическом дереве имеет не более двух потомков. Если высота дерева не превосходит $m$, то число листьев не превосходит $2^m=p$.

Но у дерева для $\omega$ не меньше $p$ листьев, значит, на некотором пути от корня к листу встречается больше $m$ нетерминальных вершин.

Поскольку различных нетерминалов всего $m$, на этом пути некоторый нетерминал повторяется. Обозначим его $A$.

Тогда в дереве есть два вхождения $A$, одно выше другого. Это даёт выводы:
\[
S\derivesstar uAz,
\]
\[
A\derivesstar vAy,
\]
\[
A\derivesstar x.
\]

Следовательно:
\[
S\derivesstar uAz
\derivesstar uvAyz
\derivesstar uvxyz.
\]

Так как $A\derivesstar vAy$, этот фрагмент можно повторить любое число раз:
\[
A\derivesstar vAy\derivesstar vvAyy\derivesstar \cdots \derivesstar v^i A y^i.
\]

После замены последнего $A$ на $x$ получаем:
\[
uv^i x y^i z\in L
\]
для всех $i\ge 0$.

Условие $|vy|\ge 1$ выполняется, потому что иначе повторение нетерминала не увеличивало бы дерево, а в нормальной форме Хомского бесконечные пустые циклы отсутствуют.

Условие $|vxy|\le p$ обеспечивается выбором нижней пары повторяющихся нетерминалов на пути: соответствующее поддерево имеет высоту не больше $m$, значит, его выход содержит не больше $2^m=p$ терминалов.

Лемма доказана.
\end{proof}

\begin{example}
Докажем, что
\[
L=\{a^n b^n c^n\mid n\ge 1\}
\]
не является контекстно-свободным.

Предположим, что $L$ контекстно-свободный. Тогда существует число накачки $p$.

Возьмём цепочку:
\[
\omega=a^p b^p c^p.
\]

По лемме:
\[
\omega=uvxyz,
\]
\[
|vxy|\le p,\qquad |vy|\ge 1.
\]

Так как $|vxy|\le p$, подцепочка $vxy$ не может одновременно содержать буквы $a$, $b$ и $c$, потому что блоки $a^p$, $b^p$, $c^p$ имеют длину $p$.

Значит, накачка $v$ и $y$ изменяет количество не всех трёх типов символов.

При $i=0$ получаем:
\[
uxz.
\]
В этой цепочке хотя бы один из трёх блоков стал короче, а хотя бы один другой остался прежней длины.

Следовательно, количества $a$, $b$ и $c$ уже не равны. Значит:
\[
uxz\notin L.
\]

Противоречие с леммой о накачке. Поэтому язык не является контекстно-свободным.
\end{example}

\section{Свойства контекстно-свободных языков}

Класс КС-языков замкнут относительно:
\begin{itemize}
    \item объединения;
    \item конкатенации;
    \item итерации Клини;
    \item гомоморфизма;
    \item обратного гомоморфизма;
    \item пересечения с регулярным языком.
\end{itemize}

Класс КС-языков не замкнут относительно:
\begin{itemize}
    \item пересечения двух КС-языков;
    \item дополнения;
    \item разности.
\end{itemize}

Для КС-языков разрешимы:
\begin{itemize}
    \item принадлежность цепочки языку;
    \item пустота языка;
    \item конечность языка.
\end{itemize}

Для КС-языков неразрешимы:
\begin{itemize}
    \item эквивалентность двух произвольных КС-грамматик;
    \item неоднозначность произвольной КС-грамматики;
    \item универсальность, то есть вопрос $L(G)=\Sigma^*$.
\end{itemize}

\section{Использование формальных языков в лексическом и синтаксическом анализе}

\subsection{Общая схема транслятора}

В компиляторах и интерпретаторах формальные языки используются для описания структуры программ.

Обычно выделяют:
\begin{enumerate}
    \item \textbf{лексический анализ};
    \item \textbf{синтаксический анализ};
    \item \textbf{семантический анализ};
    \item генерацию промежуточного представления;
    \item оптимизацию;
    \item генерацию кода.
\end{enumerate}

Формальные языки особенно важны на первых двух этапах.

\subsection{Лексический анализ}

\begin{definition}
\textbf{Лексический анализ} --- этап обработки программы, на котором входной поток символов разбивается на токены.
\end{definition}

Токенами являются:
\begin{itemize}
    \item идентификаторы;
    \item ключевые слова;
    \item числовые литералы;
    \item строковые литералы;
    \item операторы;
    \item разделители.
\end{itemize}

Лексические структуры языков программирования обычно являются регулярными языками.

\begin{example}
Идентификатор можно задать регулярным выражением:
\[
(letter)(letter|digit)^*.
\]

Целое число со знаком:
\[
(+|-)?digit^+.
\]
\end{example}

\subsubsection*{Типичная схема построения лексического анализатора}

\begin{enumerate}
    \item Для каждого класса токенов задаётся регулярное выражение.
    \item Регулярные выражения объединяются в одно большое выражение.
    \item По регулярному выражению строится НКА.
    \item НКА преобразуется в ДКА.
    \item ДКА минимизируется.
    \item Полученный ДКА используется как сканер.
\end{enumerate}

\subsubsection*{Минимальный пример}

Пусть нужно распознавать:
\[
ID=letter(letter|digit)^*,
\]
\[
NUM=digit^+.
\]

Тогда лексер работает так:
\begin{enumerate}
    \item читает символ;
    \item переходит по состояниям ДКА;
    \item запоминает последнее допускающее состояние;
    \item выбирает самый длинный подходящий префикс;
    \item возвращает соответствующий токен.
\end{enumerate}

Например, вход:
\[
abc123+45
\]
разбивается на:
\[
ID(abc123),\quad PLUS,\quad NUM(45).
\]

\subsection{Синтаксический анализ}

\begin{definition}
\textbf{Синтаксический анализ} --- этап, на котором по потоку токенов строится синтаксическая структура программы, обычно в виде дерева разбора или абстрактного синтаксического дерева.
\end{definition}

Синтаксис языков программирования обычно описывается контекстно-свободными грамматиками.

\begin{example}
Грамматика арифметических выражений:
\[
E\to E+T\mid T,
\]
\[
T\to T*F\mid F,
\]
\[
F\to (E)\mid id.
\]
\end{example}

Эта грамматика задаёт приоритет операций: умножение имеет больший приоритет, чем сложение.

\subsection{Нисходящий синтаксический анализ}

Нисходящий анализ строит дерево разбора от корня к листьям.

Наиболее простой вариант --- \textbf{метод рекурсивного спуска}.

\subsubsection*{Идея}

Для каждого нетерминала создаётся процедура.

Если есть правило:
\[
A\to BCd,
\]
то процедура $A$ вызывает процедуры $B$, $C$, а затем проверяет терминал $d$.

\begin{example}
Для грамматики:
\[
S\to aSb\mid c
\]
процедура $S$ логически выглядит так:
\begin{verbatim}
S():
    if current == 'a':
        match('a')
        S()
        match('b')
    else if current == 'c':
        match('c')
    else:
        error()
\end{verbatim}
\end{example}

\subsubsection*{Ограничения}

Рекурсивный спуск без возвратов плохо работает с:
\begin{itemize}
    \item левой рекурсией;
    \item общими префиксами альтернатив;
    \item неоднозначными грамматиками.
\end{itemize}

Левую рекурсию вида
\[
A\to A\alpha\mid \beta
\]
заменяют на:
\[
A\to \beta A',
\]
\[
A'\to \alpha A'\mid \eps.
\]

Общие префиксы устраняют левой факторизацией:
\[
A\to \alpha\beta_1\mid \alpha\beta_2
\]
заменяется на:
\[
A\to \alpha A',
\]
\[
A'\to \beta_1\mid \beta_2.
\]

\subsection{Множества FIRST и FOLLOW}

\begin{definition}
\[
FIRST(\alpha)
\]
--- множество терминалов, с которых могут начинаться цепочки, выводимые из $\alpha$.
\end{definition}

\begin{definition}
\[
FOLLOW(A)
\]
--- множество терминалов, которые могут стоять непосредственно справа от нетерминала $A$ в некоторой сентенциальной форме.
\end{definition}

Эти множества используются при построении LL-анализаторов.

\subsubsection*{Алгоритм построения FIRST}

\begin{enumerate}
    \item Если $a$ терминал, то:
    \[
    FIRST(a)=\{a\}.
    \]
    \item Если есть правило $A\to\eps$, то добавить $\eps$ в $FIRST(A)$.
    \item Если есть правило
    \[
    A\to X_1X_2\ldots X_k,
    \]
    добавить в $FIRST(A)$ все терминалы из $FIRST(X_1)$, кроме $\eps$.
    \item Если $X_1$ nullable, добавить $FIRST(X_2)$, и так далее.
    \item Повторять до стабилизации.
\end{enumerate}

\subsubsection*{Алгоритм построения FOLLOW}

\begin{enumerate}
    \item Добавить конец входа $\$$ в $FOLLOW(S)$.
    \item Для каждого правила
    \[
    A\to \alpha B\beta
    \]
    добавить
    \[
    FIRST(\beta)\setminus\{\eps\}
    \]
    в $FOLLOW(B)$.
    \item Если $\beta\derivesstar\eps$, добавить $FOLLOW(A)$ в $FOLLOW(B)$.
    \item Повторять до стабилизации.
\end{enumerate}

\subsection{Восходящий синтаксический анализ}

Восходящий анализ строит дерево от листьев к корню.

Основная операция --- \textbf{сдвиг-свёртка}.

\begin{itemize}
    \item \textbf{Сдвиг}: очередной входной символ переносится в стек.
    \item \textbf{Свёртка}: верхушка стека, совпадающая с правой частью правила, заменяется левой частью.
\end{itemize}

\begin{example}
Пусть грамматика:
\[
S\to aSb\mid ab.
\]
Разбор цепочки $aabb$ снизу вверх:

\[
aabb
\]
Сначала сворачиваем среднее $ab$:
\[
aSb
\]
Затем сворачиваем всё по правилу $S\to aSb$:
\[
S.
\]
Цепочка допущена.
\end{example}

\subsection{Алгоритм CYK}

Алгоритм Кока--Янгера--Касами проверяет принадлежность цепочки КС-языку.

Требование: грамматика должна быть в нормальной форме Хомского.

\subsubsection*{Вход}

\[
G=(V_T,V_N,P,S)
\]
в НФХ и цепочка:
\[
\omega=a_1a_2\ldots a_n.
\]

\subsubsection*{Идея}

Строится таблица $T[i,j]$, где:
\[
T[i,j]
\]
--- множество нетерминалов, из которых выводится подцепочка длины $i$, начинающаяся в позиции $j$.

\subsubsection*{Шаги алгоритма}

\begin{enumerate}
    \item Для всех $j$ заполнить:
    \[
    T[1,j]=\{A\mid A\to a_j\in P\}.
    \]

    \item Для длины $i=2,\ldots,n$:
    \[
    T[i,j]=
    \{A\mid A\to BC,\ B\in T[k,j],\ C\in T[i-k,j+k]\}
    \]
    для всех разбиений $k=1,\ldots,i-1$.

    \item Цепочка принадлежит языку тогда и только тогда, когда:
    \[
    S\in T[n,1].
    \]
\end{enumerate}

\subsubsection*{Сложность}

Временная сложность:
\[
O(n^3).
\]

\section{Краткое сравнение ролей классов языков}

\begin{center}
\begin{tabular}{>{\centering\arraybackslash}m{3cm}
                >{\centering\arraybackslash}m{4cm}
                >{\centering\arraybackslash}m{5cm}}
\toprule
Класс & Формальное средство & Использование в компиляторах \\
\midrule
Регулярные языки & Регулярные выражения, конечные автоматы & Лексический анализ \\
Контекстно-свободные языки & КС-грамматики, МП-автоматы & Синтаксический анализ \\
Контекстно-зависимые языки & КЗ-грамматики & Семантические ограничения, типы, согласование имён \\
Тип 0 & Машины Тьюринга & Общая вычислимость \\
\bottomrule
\end{tabular}
\end{center}

\section{Итоговые тезисы}

\begin{enumerate}
    \item Формальный язык --- множество цепочек над конечным алфавитом.
    \item Языки можно задавать грамматиками, автоматами, регулярными выражениями и множественно-теоретическими условиями.
    \item Иерархия Хомского классифицирует грамматики по форме правил.
    \item Регулярные языки эквивалентны конечным автоматам и регулярным выражениям.
    \item Контекстно-свободные языки эквивалентны недетерминированным автоматам с магазинной памятью.
    \item Регулярные языки применяются в лексическом анализе.
    \item Контекстно-свободные грамматики применяются в синтаксическом анализе.
    \item Для регулярных языков разрешимы многие важные задачи, включая эквивалентность.
    \item Для КС-языков задача принадлежности разрешима, но эквивалентность и неоднозначность в общем случае неразрешимы.
    \item Для практического синтаксического анализа используются специальные подклассы КС-грамматик: LL, LR, грамматики предшествования и другие.
\end{enumerate}

\end{document}
